Which is more luxurious: a diamond ring or a flower found on a path? This question is not rhetorical. This paper develops a semiotics of luxury in intimate relations organised around three doctrines: the doctrine of nature and art (luxury as the transformation of natural materials through human craft); the doctrine of sign exchange value (following Baudrillard's critique of the consumer society, in which luxury functions as a system of social differentiation); and the doctrine of relational luxury, the paper's central concept.

\emb{Relational luxury is defined not by economic value but by the degree of coupling between an object and a relational subject's spatiotemporal structure}---the extent to which a subject's existence, their intentionality and presence and choice made specifically for this relational other, is embedded in the object. A flower found on a path may be more luxurious than any diamond; a diamond purchased by proxy may be relationally impoverished despite its economic magnificence. The coupling degree is developed along four dimensions (temporal depth, intentional intensity, non-substitutability, existential investment), producing a fourfold classification in which economic and relational luxury appear as orthogonal axes.

The paper then develops a justice theory of choice. The problem of \emb{subject-embedding misrecognition} arises when an object is presented as carrying high subject-embedding but does not. The problem of \emb{equivalent agency} arises when a subject who cannot provide existential embedding $A$ substitutes another form $B$: such substitution is structurally heterogeneous, but may be just in the generative sense if four \emph{necessary but not sufficient} conditions hold---symmetric generability, real constraint, genuine embedding, and transparency. Finally, selection justice holds that the choice of relational over economic luxury is a justice stance: a refusal to use exchangeable symbols to bear what is inechangeable.

This is \pinkref{Paper~XVI} of the \emph{Philosophy of Intimacy and the Theory of Justice} series, extending the political economy of \pinkref{Paper~XV} \parencite{paperxv}.
